% !TEX root = ../main.tex
\corrigesChapitre{22}{Diffraction et interférences}
\corrigesSouples

\begin{blocLivre}{cor-22-1}
\begin{corrige}{1}{Diffraction ou pas ?}
\demarche{La diffraction est marquée lorsque la taille de l'ouverture est voisine de la longueur d'onde, ou plus petite.}
\begin{questions}
\item La porte, large d'environ \qty{1}{m}, est comparable aux longueurs d'onde sonores : le son est \res{\text{fortement diffracté}}
et se répand dans la pièce voisine. Pour la lumière, $a\gg\lambda$ : elle se propage pratiquement en ligne droite.
\item Lorsque la taille $a$ de l'ouverture est \res{\text{voisine ou inférieure à }\lambda}.
\item \res{\text{Non}} : la fréquence et la longueur d'onde sont inchangées ; seule la direction de propagation est modifiée.
\item Le son entendu derrière un obstacle, la houle qui contourne une jetée, le halo d'un lampadaire vu à travers un
voilage\,\dots
\end{questions}
\end{corrige}
\end{blocLivre}
\begin{blocLivre}{cor-22-2}
\begin{corrige}{2}{Mesurer la largeur d'une fente}
\demarche{Petits angles : $\tan\theta\approx\theta=\dfrac{L/2}{D}$, puis $a=\lambda/\theta$.}
\begin{questions}
\item La fente, l'écran à la distance $D$, la tache de largeur $L$ et le demi-angle $\theta$ entre l'axe et le bord de la tache.
\item $\theta=\dfrac{L}{2D}=\dfrac{0{,}023}{3{,}60}=\res{\qty{6,4e-3}{rad}}$.
\item $a=\dfrac{\lambda}{\theta}=\res{\qty{0,10}{mm}}$.
\item $\theta$ est inversement proportionnel à $a$ : la tache centrale devient \res{\text{deux fois plus large}}.
\end{questions}
\schema[corrige]{rev-22-ex02-cor}
\legende{Le demi-angle correspond à la moitié de la largeur de la tache centrale. Schéma sans échelle.}
\end{corrige}
\end{blocLivre}
\begin{blocLivre}{cor-22-3}
\begin{corrige}{3}{La houle entre dans le port}
\demarche{$\theta\approx\lambda/a$ ; diffraction marquée si $a$ est du même ordre que $\lambda$.}
\begin{questions}
\item \res{\text{Oui}} : l'entrée ($a=\qty{60}{m}$) est du même ordre de grandeur que la longueur d'onde (\qty{40}{m}).
\item $\theta=\dfrac{40}{60}=\res{\qty{0,67}{rad}}$, soit environ \res{38^\circ}.
\item La houle diffractée s'étale largement dans le port et atteint des bateaux \res{\text{hors de l'axe}} de l'entrée.
\item $\theta=\dfrac{40}{400}=\qty{0,10}{rad}$, soit $5{,}7^\circ$ : la houle resterait \res{\text{concentrée}} dans l'axe de la passe.
\end{questions}
\end{corrige}
\end{blocLivre}
\begin{blocLivre}{cor-22-4}
\begin{corrige}{4}{Constructives ou destructives ?}
\demarche{$\delta=S_2M-S_1M$ ; $\delta/\lambda$ entier : constructives ; demi-entier : destructives.}
\begin{questions}
\item Constructives : $\delta=k\,\lambda$ ; destructives : $\delta=\left(k+\frac12\right)\lambda$.
\item $M$ : $\delta=\qty{3,0}{cm}$, $\delta/\lambda=1{,}5$ ; $N$ : \qty{4,0}{cm}, 2 ; $P$ : 0, 0 ; $Q$ : \qty{5,0}{cm}, 2,5.
\item Eau immobile en \res{M} et en \res{Q} (demi-entiers) ; agitation maximale en \res{N} et en \res{P} (entiers).
\item Sur la \res{\text{médiatrice}} de $[S_1S_2]$ : $\delta=0$, les interférences y sont constructives et l'amplitude est
maximale.
\end{questions}
\end{corrige}
\end{blocLivre}
\begin{blocLivre}{cor-22-5}
\begin{corrige}{5}{Les fentes d'Young}
\demarche{$i=\dfrac{\lambda D}{b}$, proportionnel à $\lambda$ et à $D$, inversement proportionnel à $b$.}
\begin{questions}
\item $i=\dfrac{\num{532e-9}\times2{,}00}{\num{0.25e-3}}=\res{\qty{4,26}{mm}}$.
\item $i$ est proportionnel à $D$ : \res{\qty{8,51}{mm}}.
\item $i=\dfrac{\num{650e-9}\times2{,}00}{\num{0.25e-3}}=\res{\qty{5,20}{mm}}$.
\item $i$ est inversement proportionnel à $b$ : les franges \res{\text{s'écartent}}.
\end{questions}
\end{corrige}
\end{blocLivre}
\begin{blocLivre}{cor-22-6}
\begin{corrige}{6}{Mesurer un cheveu}
\demarche{$\theta=\dfrac{L}{2D}$ et $a=\dfrac{\lambda}{\theta}$.}
\begin{questions}
\item $\theta=\dfrac{0{,}042}{4{,}00}=\res{\qty{1,05e-2}{rad}}$.
\item $a=\dfrac{\num{633e-9}}{\num{1.05e-2}}=\res{\qty{60}{\micro m}}$.
\item Oui : la valeur est \res{\text{plausible}}.
\item Une règle graduée au millimètre ne peut pas mesurer \qty{0,06}{mm} ; la diffraction transforme cette dimension
minuscule en une tache de \res{\text{plusieurs centimètres}}, facile à mesurer.
\end{questions}
\end{corrige}
\end{blocLivre}
\begin{blocLivre}{cor-22-7}
\begin{corrige}{7}{Diffraction de la lumière blanche}
\demarche{$L=\dfrac{2\lambda D}{a}$, proportionnelle à $\lambda$.}
\begin{questions}
\item Bleu : $L=\res{\qty{2,4}{cm}}$ ; rouge : $L=\res{\qty{4,5}{cm}}$.
\item Le \res{\text{rouge}}, de plus grande longueur d'onde.
\item Le centre, éclairé par toutes les couleurs, est \res{\text{blanc}} ; les bords sont \res{\text{irisés}}, rouges vers
l'extérieur.
\item Un laser est \res{\text{monochromatique}} : les taches de diffraction sont nettes et leur largeur correspond à une seule
longueur d'onde.
\end{questions}
\end{corrige}
\end{blocLivre}
\begin{blocLivre}{cor-22-8}
\begin{corrige}{8}{La somme de deux signaux en Python}
\demarche{L'amplitude est la valeur maximale atteinte par la somme sur une période.}
\begin{questions}
\item \code{instants} contient mille dates régulièrement réparties sur une période ; \code{s} contient les valeurs
de la \res{\text{somme des deux signaux}} à ces dates.
\item \code{max(s)} donne la plus grande valeur de la somme, c'est-à-dire son \res{\text{amplitude}}. La somme est périodique,
de même période : une période suffit.
\item $\varphi=0$ : $2A=2$ ; $\varphi=\pi/2$ : $2\cos(\pi/4)=1{,}41$ ; $\varphi=2\pi/3$ : $2\cos(\pi/3)=1$ ; $\varphi=\pi$ : 0. Les
valeurs affichées sont \res{\text{conformes}}.
\item On crée la liste \code{s1} des valeurs de $\sin(2\pi t)$ et la liste \code{s2} des valeurs de $\sin(2\pi t+\pi/2)$, puis
leur somme : \code{somme = [x + y for x, y in zip(s1, s2)]}. On trace chaque liste en fonction de \code{instants} avec
\code{plt.plot}, puis on affiche le graphique avec \code{plt.show()}.
\end{questions}
\schema[corrige]{rev-22-ex08-cor}
\legende{Deux signaux de même fréquence déphasés de $\pi/2$ ; l’amplitude de la somme vaut $\sqrt2$ unité.}
\end{corrige}
\end{blocLivre}
\begin{blocLivre}{cor-22-9}
\begin{corrige}{9}{Mesurer une longueur d'onde}
\demarche{$i=\dfrac{\qty{3,8}{cm}}{10}$, puis $\lambda=\dfrac{i\,b}{D}$.}
\begin{questions}
\item L'incertitude de lecture porte sur la distance totale : elle est \res{\text{divisée par 10}} pour un interfrange.
\item $i=\res{\qty{3,8}{mm}}$ et $\lambda=\dfrac{\num{3.8e-3}\times\num{0.40e-3}}{2{,}50}=\res{\qty{608}{nm}}$ : lumière \res{\text{orangée}}.
\item $\dfrac{\qty{1}{mm}}{\qty{38}{mm}}\approx\res{\qty{2,6}{\percent}}$.
\item Franges brillantes : $\dfrac{b\,x}{D}=k\lambda$, soit $x_k=k\,\dfrac{\lambda D}{b}$ ; deux franges consécutives sont séparées de
$x_{k+1}-x_k=\res{\dfrac{\lambda D}{b}}$.
\end{questions}
\end{corrige}
\end{blocLivre}
\begin{blocLivre}{cor-22-10}
\begin{corrige}{10}{La cuve à ondes}
\demarche{$\lambda=v/f$ ; points de repos : $\delta=(k+\frac12)\lambda$ ; maximums : $\delta=k\lambda$, avec $|\delta|<\qty{6,0}{cm}$.}
\begin{questions}
\item $\lambda=\dfrac{0{,}30}{20}=\res{\qty{1,5}{cm}}$.
\item Repos : $|k+\frac12|\times1{,}5<6{,}0$, soit $k$ de $-4$ à 3 : \res{\text{8 lignes de repos}}. Maximums : $|k|\times1{,}5<6{,}0$, soit
$k$ de $-3$ à 3 : \res{\text{7 lignes}}.
\item Sur le segment, un déplacement de $x$ change $\delta$ de $2x$ ; $\delta$ augmente de $\lambda$ entre deux points de repos
consécutifs, d'où un écart de \res{\lambda/2=\qty{0,75}{cm}}.
\item $\lambda$ serait divisée par 2 : deux fois plus de lignes, \res{\text{deux fois plus serrées}}.
\end{questions}
\end{corrige}
\end{blocLivre}
\begin{blocLivre}{cor-22-11}
\begin{corrige}{11}{Deux sources incohérentes}
\demarche{Des interférences stables exigent deux ondes de même fréquence et de déphasage constant.}
\begin{questions}
\item \res{\text{Non}} : deux lampes émettent des trains d'ondes sans relation de phase ; le déphasage varie au hasard et
les franges se brouillent.
\item Les deux fentes sont éclairées par la \res{\text{même onde}} : les ondes qui en sont issues gardent un déphasage
constant.
\item Les haut-parleurs alimentés par le même générateur, \res{\text{oui}} ; deux musiciens, non : leurs fréquences ne sont
jamais rigoureusement égales et leur déphasage varie.
\item Au centre, $\delta=0$ pour toutes les couleurs : la frange est \res{\text{blanche}}. Ailleurs, l'interfrange $i=\lambda D/b$
dépend de la couleur : les franges des différentes couleurs se décalent et se mélangent en irisations.
\end{questions}
\end{corrige}
\end{blocLivre}
\begin{blocLivre}{cor-22-12}
\begin{corrige}{12}{La finesse des images d'un télescope}
\demarche{$\theta=1{,}22\,\lambda/D$ ; taille du plus petit détail $=\theta\times$ distance.}
\begin{questions}
\item $\theta=1{,}22\times\dfrac{\num{550e-9}}{2{,}4}=\res{\qty{2,8e-7}{rad}}$.
\item $\num{2.8e-7}\times\num{3.84e8}=\res{\qty{107}{m}}$.
\item Œil : $\theta=\qty{2,2e-4}{rad}$, soit des détails de \qty{86}{km} sur la Lune : Hubble voit des détails \res{\text{800 fois}}
plus petits.
\item À \qty{2}{\micro m}, $\theta=\qty{3,8e-7}{rad}$ : la résolution du James-Webb est \res{\text{comparable}} à celle de Hubble, car il
observe à plus grande longueur d'onde. Un grand miroir collecte plus de lumière et, à longueur d'onde égale, distingue
des détails plus fins.
\end{questions}
\end{corrige}
\end{blocLivre}
\begin{blocLivre}{cor-22-13}
\begin{corrige}{13}{Le casque à réduction de bruit active}
\demarche{Opposition de phase : retard d'une demi-période ; amplitude résiduelle $2A\sin(\pi f\Delta t)\leqslant A/10$.}
\begin{questions}
\item Une demi-période : $\dfrac{T}{2}=\dfrac{1}{2\times200}=\res{\qty{2,5}{ms}}$.
\item Pour $\Delta t=0$, $\sin0=0$ : les deux ondes s'annulent \res{\text{exactement}}.
\item $2\pi f\,\Delta t\leqslant0{,}1$, soit $\Delta t\leqslant\dfrac{0{,}05}{\pi f}$ : \res{\qty{80}{\micro s}} à \qty{200}{Hz},
\res{\qty{8,0}{\micro s}} à \qty{2000}{Hz}.
\item Le grondement d'un moteur est grave et régulier : l'électronique a le temps de le compenser. Les voix contiennent
des fréquences élevées, qui exigent un retard minuscule, et varient sans cesse : elles sont \res{\text{mal atténuées}}.
\end{questions}
\end{corrige}
\end{blocLivre}
\begin{blocLivre}{cor-22-14}
\begin{corrige}{14}{Le traitement antireflet}
\demarche{Réflexion supprimée si $\delta=2\,n\,e=\left(k+\frac12\right)\lambda$ ; plus petite épaisseur pour $k=0$.}
\begin{questions}
\item Interférences destructives : $\res{2\,n\,e=\left(k+\frac12\right)\lambda}$.
\item $k=0$ : $e=\dfrac{\lambda}{4n}=\dfrac{550}{4\times1{,}38}=\res{\qty{99,6}{nm}}$, environ \qty{100}{nm}.
\item Bleu : $\delta/\lambda=\res{0{,}61}$ ; rouge : $\delta/\lambda=\res{0{,}39}$. Ces valeurs, éloignées de 0,5, ne donnent qu'une
extinction \res{\text{partielle}}.
\item Le vert, au centre du spectre, est éteint ; une partie du bleu et du rouge reste réfléchie : leur mélange donne
une teinte \res{\text{violacée}}.
\end{questions}
\end{corrige}
\end{blocLivre}
\begin{blocLivre}{cor-22-15}
\begin{corrige}{15}{Mesurer la vitesse du son par interférences}
\demarche{Le micro s'éloigne d'une source et se rapproche de l'autre : $\delta$ varie de $2x$.}
\begin{questions}
\item Une distance augmente de $x$, l'autre diminue de $x$ : $\delta$ varie de \res{2x}.
\item Entre deux minimums successifs, $\delta$ varie de $\lambda$ : le micro s'est déplacé de \res{\lambda/2}.
\item Dix minimums délimitent 9 intervalles : $\dfrac{\lambda}{2}=\dfrac{77{,}4}{9}=\qty{8,6}{cm}$, d'où $\lambda=\res{\qty{17,2}{cm}}$ et
$v=\lambda f=\res{\qty{344}{m.s^{-1}}}$.
\item L'erreur de repérage d'un minimum se répartit sur neuf intervalles : l'\res{\text{incertitude relative}} est
nettement réduite.
\end{questions}
\end{corrige}
\end{blocLivre}
